\chapter{Fondamenti}\label{chap:background}
\section{Modello autoregressivo}
Dato un prompt $x$ e una glossa tokenizzata $y$, il modello fattorizza
$p_\theta(y\mid x)=\prod_t p_\theta(y_t\mid x,y_{<t})$. Il tokenizer opera su
subtoken, mentre metriche e vocabolario di glosse operano soprattutto su token
separati da spazi: confondere i due livelli produce gate e vincoli errati.

\section{Supervisione, policy optimization e vincoli}
SFT apprende dai target gold; GRPO confronta più completion dello stesso prompt;
il constraint modifica il supporto disponibile in generazione. Sono interventi
diversi: un Trie può garantire compatibilità lessicale senza migliorare il
contenuto, e un reward ben formulato può essere empiricamente poco informativo.
BLEU \cite{papineni2002bleu}, ROUGE \cite{lin2004rouge} e chrF misurano aspetti
diversi e vanno letti insieme.

\section{Riproducibilità}
Un numero senza identità di dataset, modello, tokenizer, vocabolario, retrieval,
grammar, decoding, seed e budget non è un risultato riproducibile. Il progetto
materializza queste identità negli artifact e usa ID stabili per cache e
bootstrap appaiato. Provenienza e failure espliciti sono parte del metodo.
